%% kapa: exam=1 round=34 year=2023 period=2 form=A subject=appraisal_law subject_ko=감정평가관계법규 seq=15 number=15
%% source: Q-Net 감정평가사 기출문제 (cst00309) artlSeq=5211789 file=raw/1008/5211789/2023년도 제34회 감정평가사 1차시험 2교시.pdf
%% conversion: kordoc 4.24.51 -> tools/gen1.py (auto); review=proofread-sonnet; flags=none
\begin{question}{15}
부동산 가격공시에 관한 법령상 개별공시지가에 관한 설명으로 옳지 않은 것은?
\begin{choices}
  \item 시·도지사는 개별공시지가를 산정한 때에는 중앙부동산가격공시위원회의 심의를 거쳐 이를 공시하여야 한다.
  \item 표준지로 선정된 토지에 대한 개별공시지가가 결정·공시되지 아니한 경우 해당 토지의 표준지공시지가를 개별공시지가로 본다.
  \item 공시기준일 이후에 토지의 형질변경으로 「공간정보의 구축 및 관리 등에 관한 법률」에 따른 지목변경이 된 토지는 개별공시지가 공시기준일을 다르게 할 수 있는 토지에 해당한다.
  \item 개별공시지가의 결정·공시에 소요되는 비용 중 국고에서 보조할 수 있는 비용은 개별공시지가의 결정·공시에 드는 비용의 50퍼센트 이내로 한다.
  \item 토지가격비준표의 적용에 오류가 있는 경우는 개별공시지가를 정정하여야 할 사유에 해당한다.
\end{choices}
\end{question}
\answer{1}
